\documentclass[conference]{IEEEtran}

\usepackage{graphicx}
\usepackage{cite}
\usepackage{amsmath,amssymb}
\usepackage{booktabs}
\usepackage{placeins} % provides \FloatBarrier to stop figure drift across sections

% Figures are stored in ./figures
\graphicspath{{figures/}}

\title{A Novel Approach to Secure Execution in Embedded Systems: Custom ISA and Compiler Binding}
\author{mariaayub}

\begin{document}
\maketitle

\section{Abstract}
Embedded devices increasingly operate in adversarial environments where physical access
enables firmware replacement and invasive debugging. Although secure boot protects the
initial image, it does not necessarily constrain post-boot debug and programming
interfaces. This paper presents a hardware--software co-design that gates both execution
and debug enablement on a device-resident authorization decision. The approach adds a
minimal ISA extension with a custom authorization instruction (\texttt{kc}) that invokes
verification against key material stored in protected non-volatile memory, never exposed
via the standard load/store path. A constant-time comparator and security FSM enforce
fail-stop behavior: execution and debug are released only after successful verification.
On the host side, authorization material is derived from host and device identifiers and
recorded in a binding database to support traceability and audit. Toolchain integration
ensures the authorization primitive is emitted and preserved across compilation and
image generation. Evaluation artifacts show end-to-end feasibility, including successful
emission of the custom instruction, key exchange preceding debug enablement, and
workflow logs consistent with the intended policy. Index Terms---embedded systems,
debug access control, ISA extensions, hardware--software co-design, key-based
authorization, toolchain integration.

\section{Introduction}
Embedded systems now underpin critical infrastructure, industrial control, and
consumer devices, yet many remain vulnerable to unauthorized reflashing and invasive
debugging when an attacker gains physical access. Traditional secure boot mechanisms
provide an important integrity check at startup, but they do not necessarily restrict
post-boot access to programming and debug capabilities. In practice, debug interfaces
remain a powerful attack surface because they expose memory and execution control that
can bypass higher-level software defenses.

This work proposes a device-resident authorization mechanism that binds privileged
capabilities to an explicit, ISA-visible trigger. The core idea is to couple an
authorization instruction to a protected verification datapath so that authorization
is invoked as part of the execution model, not merely as a host-side convention. The
device stores a reference verifier in protected non-volatile memory that is not
memory-mapped into the normal address space; a constant-time compare unit and a small
security FSM decide whether to release \texttt{exec\_enable} and \texttt{debug\_enable}.
The host tooling derives authorization material from host and device identifiers,
records the binding, and communicates the authorization sequence prior to enabling
privileged functionality.

The contributions of this paper are threefold: (i) a minimal ISA extension and hardware
control path that trigger authorization without exposing key material, (ii) a binding
model and host workflow that associate authorization material with specific devices and
operators, and (iii) an end-to-end toolchain integration and evidence-based evaluation
showing that the custom instruction is emitted in firmware images and that debug access
is gated by successful authorization. Collectively, these contributions provide a
deployable path for instruction-triggered, key-gated authorization in embedded systems.

The remainder of the paper is organized as follows. Section~\ref{sec:relatedwork}
reviews related work in secure boot, debug protection, and hardware-enforced security.
Section~\ref{sec:architecture} details the system architecture and authorization
dataflow. Section~\ref{sec:results} presents the evaluation artifacts and functional
validation. Section~\ref{sec:discussion} discusses security implications and operational
considerations, Section~\ref{sec:futurework} outlines future directions, and
Section~\ref{sec:conclusion} concludes.

\input{paper/RELATED_WORK.tex}

\input{paper/ARCHITECTURE.tex}

\input{paper/RESULTS.tex}

\input{paper/DISCUSSION.tex}

\input{paper/FUTURE_WORK.tex}

% Paper-ready Conclusion section (IEEE style).

\section{Conclusion}
\label{sec:conclusion}

This work presented a hardware--software co-design for securing embedded programming
and execution under physical-access threat models. The architecture gates normal
execution and invasive debug access on a device-resident authorization decision, using
protected key storage and a verification datapath that is not exposed via ordinary
memory transactions. Authorization is invoked through an ISA-visible primitive
integrated into the toolchain, enabling reproducible and inspectable binaries.

Evaluation artifacts demonstrate end-to-end feasibility: the custom instruction is
compiled and preserved in generated images, key exchange and matching precede debug
enablement, the host workflow supports compilation and programming, and failure
results in fail-stop behavior. Together, these results support the practicality of
instruction-triggered, key-gated control for restricting unauthorized reflashing and
debugging in embedded systems.



\bibliographystyle{IEEEtran}
\bibliography{refs}
\end{document}

